% Section 1 of the distilled proof: everything the readout is made of.
% Self-containment rule: no phrase of the form ``as in the specification''
% survives here without the specification's own text printed beside it.
\section{Setup}\label{sec:setup}

\subsection{Polygons, guarded predicates, genericity}

\begin{definition}[polygon]\label{def:polygon}
Let $n\geq3$. A \emph{polygon} on $n$ vertices is a tuple
$P=(p_1,\dots,p_n)\in(\mathbb{R}^2)^n$ with $p_{i+1}\neq p_i$ for every $i$,
indices read cyclically modulo $n$ throughout. Its \emph{edges} are the closed
segments $e_i=[p_i,p_{i+1}]$ and its \emph{edge directions} are
$d_i=p_{i+1}-p_i\in\mathbb{R}^2\setminus\{0\}$. Two edges are \emph{consecutive}
if their index sets meet, and \emph{remote} otherwise. An edge is \emph{remote
to a vertex} $p_M$ when it is remote to both edges incident to $M$ --- the
stronger of the two possible readings, and the one every consumer of the phrase
uses: it excludes an edge that merely avoids $p_M$ while sharing a vertex with
one of its edges. We write
$\det(u,v)=u_1v_2-u_2v_1$.
\end{definition}

\begin{definition}[principal turns and the regular locus]\label{def:regular}
\begin{enumerate}
\item[(A)] Let $L$ be a closed polygon with corners $q_0,\dots,q_{c-1}$ and edge directions
$\delta_i=q_{i+1}-q_i$. The \emph{principal turn} at $q_i$ is the unique
$\tau_i\in(-\pi,\pi)$ with
$\cos\tau_i=\frac{\langle\delta_{i-1},\delta_i\rangle}{|\delta_{i-1}||\delta_i|}$
and $\sgn\tau_i=\sgn\det(\delta_{i-1},\delta_i)$. It exists precisely when
$\delta_i$ is not a negative multiple of $\delta_{i-1}$: if
$\det(\delta_{i-1},\delta_i)\neq0$ it is nonzero, and if $\delta_i$ is a
positive multiple of $\delta_{i-1}$ it is $0$. The excluded case
$\delta_i\in\mathbb R_{<0}\delta_{i-1}$ is the \emph{kink}, where the two
candidate values $\pm\pi$ are both excluded from the open interval.
\item[(B)] The \emph{regular locus} $\mathcal R_n\subseteq(\mathbb R^2)^n$ is the set of
$P=(p_1,\dots,p_n)$ with $d_i\neq0$ for all $i$ and $d_{i+1}\neq-\lambda d_i$
for all $i$ and all $\lambda>0$: nonzero edges and no consecutive pair that
doubles back. Equivalently, every principal turn of
clause~(A) exists.
\end{enumerate}
\end{definition}

\begin{definition}[the guarded list $\mathcal G$]\label{def:guarded}
Fix $n$. The \emph{guarded list} is a finite family of real polynomial functions
on $(\mathbb R^2)^n$, indexed once and for all by combinatorial data that does
not depend on the configuration. Write $d_i=p_{i+1}-p_i$ and, for an edge index
$e$, write $\ell_e(x)=\det(d_e,\,x-p_e)$ for the affine form vanishing on the
line of $e$, with coefficient row
$\bigl(A_e,B_e,C_e\bigr)=\bigl(-d_{e,2},\;d_{e,1},\;d_{e,2}p_{e,1}-d_{e,1}p_{e,2}\bigr)$,
so that $\ell_e(x)=A_ex_1+B_ex_2+C_e$.

\emph{Unconditional members.}
\begin{enumerate}
\item[(G1)] $\mathrm{G1}_i=\det(d_{i-1},d_i)$, for every $i\in\{1,\dots,n\}$.
\item[(G2)] $\mathrm{G2}_{e,i}=\ell_e(p_i)=\det(d_e,\,p_i-p_e)$, for every edge
index $e$ and every $i\notin\{e,e+1\}$.
\end{enumerate}

\emph{Activation.} Two remote edges $e,f$ are \emph{defined} to cross when
\[
   \mathrm{G2}_{e,f}\,\mathrm{G2}_{e,f+1}<0
   \quad\text{and}\quad
   \mathrm{G2}_{f,e}\,\mathrm{G2}_{f,e+1}<0 ,
\]
a condition on unconditional members alone. The products are \emph{strict}, and
that is a decision, not a formality: an earlier revision wrote
$\sgn\neq\sgn$, which is undecided exactly where it is consumed. At a wall
point one of the four members vanishes, and the three available readings
disagree there. The strict product makes the pair \emph{inactive} at such a
point; a three-valued $\sgn\in\{-1,0,+1\}$ would make it active, and then the
sliding bundle of the wall dictionary of
Section~\ref{sec:dictionary} would be $\{\mathrm{G2},\mathrm{G4}\}$ rather than
the singleton that dictionary's item~(iii) records;
and the shipped evaluator's test treats a zero as negative, which is asymmetric
between the two edges. The strict product is taken, and the divergence from the
evaluator is recorded here rather than left silent. What can be said about its
consequences is scoped to this document: every statement below evaluates $X_1$
at \emph{generic} polygons, where no $\mathrm{G2}$ vanishes and the two
readings agree, so no statement below depends on which reading the code takes.
What an external caller passes to that code is outside this document's claims
--- a caller that hands it a collinear configuration reaches the divergence, and
nothing here says otherwise. An earlier revision wrote that the code "never
reaches" such a configuration, which is a claim about the code's callers and not
about this document. It is the sign
condition, and not the geometric phrase ``their relative interiors meet'', that
every consumer below reads and that the shipped evaluator computes.

The two agree wherever the four members are nonzero, and there they say that
each edge has its two endpoints strictly on opposite sides of the other's line,
which for segments is exactly a transverse interior crossing. They do
\emph{not} agree everywhere in $\mathcal R_n$, and an earlier revision wrote
``exactly when'' as though they did: two collinear overlapping remote edges have
meeting relative interiors while all four members vanish, so the sign condition
fails. Such a configuration is not generic --- those vanishing members are
unconditional, hence relevant --- so the generic locus defined below and the
discriminant $\mathcal D$ of Section~\ref{sec:dictionary} are the same sets under
either reading, and no statement
below changes. What changes is that activation is now decided by a formula.

\emph{Conditional members.}
\begin{enumerate}
\item[(G5)] $\mathrm{G5}_{e,f}=\det(d_e,d_f)$, for every pair of remote edge
indices with $1\leq e<f\leq n$ as integers --- the comparison is between the
integer representatives $1,\dots,n$ of the edge indices and not a cyclic
comparison, there being no such thing; adjacency and remoteness continue to be
read modulo $n$. This fixes one \emph{ordered representative} per pair, and an
unordered pair will not do, because
$\det(d_f,d_e)=-\det(d_e,d_f)$ and the family is used where signs are read, in
the sign-change condition that defines a transversal event below; \emph{active} when $e$
and $f$ cross. Making it
conditional rather than unconditional is a fidelity point: the specification
guards the turn at a smoothing site, and a smoothing site exists only where two
edges cross. Demanding $\det(d_e,d_f)\neq0$ for a pair that does not cross
would shrink the generic locus below the specification's and exclude polygons
the laws are stated about.
\item[(G3)] $\mathrm{G3}_{e,f,g}=\det\begin{pmatrix}A_e&B_e&C_e\\
A_f&B_f&C_f\\ A_g&B_g&C_g\end{pmatrix}$, the rows being the line-coefficient
rows bound at the head of this definition, one for each of $e$, $f$ and $g$;
for every triple of pairwise remote edges $1\leq e<f<g\leq n$, ordered by their
integer representatives; \emph{active} when $e,f,g$ pairwise cross. It vanishes exactly when the
three lines share a point of the projective plane: either they are concurrent in
the plane, or all three are parallel. A parallel \emph{pair} with a transverse
third does \emph{not} make it vanish --- an earlier revision said it did, and the
configuration $p_0=(0,0)$, $p_1=(1,0)$, $p_2=(0,1)$, $p_3=(1,1)$, $p_4=(2,-1)$,
$p_5=(2,2)$ refutes it, the rows of $e_0,e_2,e_4$ being $(0,1,0)$, $(0,1,-1)$
and $(-3,0,6)$ with determinant $3$. On the active domain the three edges
pairwise cross, so no two of the lines are parallel and the vanishing is exactly
a concurrence at a finite point, which is what every consumer reads.
\item[(G4)] $\mathrm{G4}_{e;f,g}
=\det(d_f,\,p_f-p_e)\,\det(d_g,d_e)-\det(d_g,\,p_g-p_e)\,\det(d_f,d_e)$,
for every edge $e$ and every pair of edges remote to $e$ with
$1\leq f<g\leq n$ as integers --- again the comparison is between integer
representatives, and again one ordered representative per pair, since
$\mathrm{G4}_{e;g,f}=-\mathrm{G4}_{e;f,g}$; \emph{active} when $f$ and $g$ both
cross $e$.

\emph{Normalization of cyclic aliases.} Edge indices are read cyclically, so
one edge has many names: $e_0$, $e_n$ and $e_{2n}$ are the same edge. Every
comparison of indices in this document --- the $f<g$ of (G5), the $e<f<g$ of
(G3), the $f<g$ of (G4) --- is a comparison of \emph{representatives}: each
index is first reduced to the unique integer in $\{1,\dots,n\}$ naming that
edge, and the reduced integers are compared. Adjacency, remoteness and the
edges themselves are unaffected by the reduction; only the choice of ordered
representative depends on it, and only that choice carries a sign.

\emph{The two accessors.} Consumers meet a pair of edges in the order the
geometry gives, which after reduction may be the reverse of the representative
order. Let $f,g$ be distinct edges remote to $e$, presented in some order and
possibly under aliases, and let $\bar f,\bar g\in\{1,\dots,n\}$ be their
representatives. Define
\[
   \mathrm{G4}\langle e;f,g\rangle
   =\mathrm{G4}_{e;\min(\bar f,\bar g),\,\max(\bar f,\bar g)}\in\mathcal G,
   \qquad
   \epsilon(f,g)=\begin{cases}+1,&\bar f<\bar g,\\ -1,&\bar f>\bar g,\end{cases}
\]
the first a \emph{member-valued} accessor --- it returns an element of the
guarded list, the one indexed by the ordered representative pair --- and the
second an \emph{orientation sign} in $\{\pm1\}$, a constant of the index data
and not a function of the configuration. Their product
\[
   \widehat{\mathrm{G4}}_{e;f,g}
   =\epsilon(f,g)\cdot\mathrm{G4}\langle e;f,g\rangle
\]
is the \emph{oriented value}: a real polynomial function on $(\mathbb R^2)^n$,
antisymmetric in the ordered pair, equal to the member up to the orientation
sign.

The two are used in different places and the distinction is not cosmetic.
Every set of guarded predicates formed below --- in particular the zero set of
an event, defined later in this section --- is a set of \emph{members} of
$\mathcal G$, and membership there is by \emph{index}: the element is
$\mathrm{G4}\langle e;f,g\rangle$, the member indexed by the ordered
representative pair, and not the oriented value, which when $\epsilon=-1$ is
that member's negative and is indexed by nothing. The distinction is one of
indexing and not of polynomials, and the difference matters: as polynomials a
negated member can coincide with another member, since for three pairwise
remote edges $\mathrm{G4}_{e;f,g}=-\mathrm{G3}_{e,f,g}$ --- an identity of
polynomials, proved in the dictionary section where the canonical factors are
analysed --- so $-\mathrm{G4}_{e;f,g}$ is the polynomial $\mathrm{G3}_{e,f,g}$,
which \emph{is} a member. An earlier revision said the
oriented value "is not a member of $\mathcal G$ when $\epsilon=-1$", which as a
statement about polynomials is false.

Statements about vanishing, activation and relevance are statements about the
member and are insensitive to the presentation order. Statements about
\emph{sign} --- a sign change across an event, a factorization whose two sides
must agree in sign, the order of two crossings along an edge --- read the
oriented value, and there the orientation sign $\epsilon$ is carried
explicitly. Both kinds occur below, and the next display is of the first kind.

Two illustrations of the reduction, both used below. At $j=0$ the pair
$(e_{j-1},e_j)$ is $(e_{n-1},e_n)$, whose representatives are already
increasing: $\epsilon=+1$, and the oriented value is the member. At $j=1$ it is
$(e_0,e_1)$, whose representatives are $(n,1)$: $\epsilon=-1$, and the member
is $\mathrm{G4}_{e;1,n}$. An earlier revision wrote
$\mathrm{G4}_{h;e_{j-1},e_j}$ where the reduction makes that name empty and the
sign opposite, and a later one placed the wrap at $j=0$, where there is none. When active, both
$\det(d_f,d_e)$ and $\det(d_g,d_e)$ are nonzero, the crossing of $e$ with $f$
sits at the parameter $t_f=\det(d_f,p_f-p_e)/\det(d_f,d_e)$ along $e$, and, for
the ordered representative pair $\bar f<\bar g$ that indexes the member,
\[
   \mathrm{G4}_{e;\bar f,\bar g}
      =(t_{\bar f}-t_{\bar g})\det(d_{\bar f},d_e)\det(d_{\bar g},d_e),
\]
so its vanishing is exactly a tie $t_{\bar f}=t_{\bar g}$ in the order of
crossings along $e$. Presented in the other order the same identity holds for the oriented value,
both sides changing sign. Two kinds of consumer read this display, and they read
different things. Those that need only the \emph{vanishing} --- the two collision arguments
below, in the chamber-invariance proposition and in the silence lemma --- read
the member, and the presentation order is irrelevant to them. Those that
need the \emph{sign} --- the order of the two crossings along $e$, and the
bigon-bundle computation of Section~\ref{sec:dictionary} --- read the oriented
value $\widehat{\mathrm{G4}}_{e;f,g}=\epsilon(f,g)\,\mathrm{G4}\langle
e;f,g\rangle$ and carry $\epsilon$. An earlier revision said the vanishing was
"the only reading the consumers of this display take", which the sign-reading
sites contradict.
\end{enumerate}
\end{definition}

\begin{definition}[relevant members, generic polygons and chambers]\label{def:generic}
\begin{enumerate}
\item[(A)] A member of $\mathcal G$ is \emph{relevant} at a polygon $P$ if it is
unconditional, or if it is conditional and active at $P$
(Definition~\ref{def:guarded}). A polygon $P$
(Definition~\ref{def:polygon}) is \emph{generic} if every member relevant at
$P$ is nonzero at $P$; equivalently, if every unconditional member of
$\mathcal G$ is nonzero at $P$ and every conditional member that is active at
$P$ is nonzero at $P$. We write $\mathcal U_n$ for the set of generic polygons
on $n$ vertices.

\item[(B)] A \emph{chamber} is a connected component of the set $\mathcal U_n$ of generic
polygons on $n$ vertices (clause~(A)). It is \emph{not} in
general a component of the set where every member of $\mathcal G$ is nonzero:
a conditional member that is inactive may vanish, or vanish identically, at a
perfectly generic polygon, and nothing in this document asks it not to.
\end{enumerate}
\end{definition}

\begin{remark}[what genericity excludes, item by item]\label{rem:genericmeans}
$\mathrm{G1}_i\neq0$ says $d_{i-1}$ and $d_i$ are linearly independent, which
excludes \emph{both} degenerate turns at once: the \emph{flat} turn, where
$d_i$ is a positive multiple of $d_{i-1}$ and $p_i$ lies strictly between its
neighbours, and the \emph{kink}, where $d_i$ is a negative multiple of $d_{i-1}$
and the polygon reverses at $p_i$. $\mathrm{G2}_{e,i}\neq0$ says no vertex lies
on the line of a non-incident edge, hence none lies on the edge.
$\mathrm{G3}\neq0$ says no three pairwise remote edges are concurrent, so no two
crossings collide. $\mathrm{G4}\neq0$ says the crossings along each edge are
totally ordered with no ties. $\mathrm{G5}\neq0$ says two crossing edges are
never parallel, so every crossing is a transversal double point. Consequently a
generic polygon meets each of its remote edges in finitely many transversal
double points, has no triple points, and has no vertex on a non-incident edge:
it is an A'Campo divide.
\end{remark}

\begin{proposition}[active (G5) and the frozen generic locus]\label{prop:fidelity}
Fix $n\geq3$. Let $\mathcal U_n^{\flat}$ be the set of polygons at which every member of the
four families (G1)--(G4) that is relevant is nonzero --- the specification's
guarded list, which does not mention $\det(d_e,d_f)$.
\begin{enumerate}
\item[(i)] Let $P$ be a polygon on $n$ vertices with $\mathrm{G1}_i(P)\neq0$ for every $i$ and
$\mathrm{G2}_{e,i}(P)\neq0$ for every edge $e$ and every valid
$i\notin\{e,e+1\}$. Then every (G5) member active at $P$ is nonzero,
equivalently every remote pair whose relative interiors meet has nonzero direction determinant.
\item[(ii)] $\mathcal U_n=\mathcal U_n^{\flat}$.
\end{enumerate}
\end{proposition}
\begin{proof}
(i) Let $e,f$ be remote with a common interior point and suppose
$\det(d_e,d_f)=0$. Then $d_e$ and $d_f$ are proportional and the two segments,
having a common point, lie on one line and overlap in a segment of positive
length or in a point. If they overlap in a segment, an endpoint of one lies in
the other, so the corresponding $\mathrm{G2}$ vanishes, contrary to hypothesis.
If they meet in a single point, that point is interior to both by assumption,
and two collinear segments meeting in exactly one interior point of each is
impossible. So $\det(d_e,d_f)\neq0$.

(ii) $\mathcal U_n\subseteq\mathcal U_n^{\flat}$ because $\mathcal U_n$ asks for more.
Conversely let $P\in\mathcal U_n^{\flat}$; every relevant (G1), (G2), (G3) and
(G4) is nonzero there. Since (G1) and (G2) are unconditional, clause (i) gives
$\mathrm{G5}_{e,f}\neq0$ for every active pair. An inactive (G5) is not
relevant, so every member relevant at $P$ is nonzero and $P\in\mathcal U_n$.
\end{proof}

\begin{remark}[why (G5) is conditional, and what it would cost to change]
\label{rem:fidelity}
Proposition~\ref{prop:fidelity}(ii) is the fidelity statement, and it is stated
rather than remarked because it is load-bearing: it says no theorem below is
proved on a narrower domain than the specification states it for.
Definition~\ref{def:guarded} adds (G5) as a \emph{conditional} member so that
the turn at a smoothing site has a named predicate, and the proposition says
that this adds nothing.

Making (G5) unconditional would \emph{not} be harmless: the pinned bow-tie $K_0$
of Definition~\ref{def:circularstar} has two parallel remote edges that do not
meet, so it would cease to be generic and the rank-zero anchor of
Theorem~\ref{thm:main}(B) would not exist. The same bow-tie later refuted an
argument that ran general position on the unconditional zero set
(Lemma~\ref{lem:relgp}(B)); it is the configuration this whole distinction exists
to keep.
\end{remark}


\subsection{The Gauss word, supports and carriers}

Fix a \emph{diagrammatic} polygon $P$ --- the class defined below, of which
every generic polygon is a member --- and let $m=m(P)$ be the number of its
double points, labelled $1,\dots,m$. Everything in this subsection, up to and
including the residual pieces and their diagrams, is stated for that class:
genericity enters later, with the guards, and an earlier revision fixed a
generic polygon here and then claimed the wider scope downstream. Traversing $P$ once from $p_1$ in the direction of
increasing index visits each double point exactly twice; recording the labels in
the order visited gives a cyclic word of length $2m$ in which every letter occurs
twice, the \emph{Gauss word} of $P$.

\begin{definition}[diagrammatic polygon]\label{def:diagrammatic}
A polygon $P$ is \emph{diagrammatic} if its self-intersections are finitely
many, each of them an isolated transverse crossing of two edges with exactly two
preimages on the traversal circle, no two of them sharing an image or a
preimage, none of them at a corner, and no vertex of $P$ lying on an edge not
incident to it. The clause about preimages and shared images is not implied by
the others, and an earlier revision omitted it: the polygon
$((-3,0),(3,0),(0,-3),(0,3),(-2,-2),(2,2))$ has its edges $e_0,e_2,e_4$ meeting
pairwise transversally --- the three determinants are $36$, $24$, $-24$ --- at
the \emph{single} point $(0,0)$, so three strands pass through one point, and the
$2m$-letter Gauss word does not exist.

Every generic polygon is diagrammatic: the (G2) members forbid a vertex on a
remote edge, the (G5) members make each crossing transversal, and a triple point
would make a (G3) member vanish while all three pairs cross, so that member is
active, hence relevant, hence nonzero. So is the polygon at a simple positive
flat wall, whose only degeneracy is a vanishing turn with the middle vertex
\emph{inside} the segment. That last qualification matters: at an exterior
collinear vertex the two incident edges overlap and the far one carries the
middle vertex, as in $((0,0),(2,0),(1,0),(0,1))$, which is not diagrammatic ---
an earlier revision said "any polygon whose only degeneracy is a vanishing
turn".

The class is named because the Gauss word, the interlacement graph, the
independent sets, the undominated sets, the residual pieces and their diagrams
are read off a diagrammatic polygon and need nothing more; genericity is what
$X_1$ needs, through the guards, and is a strictly stronger condition. Three
sections use exactly this class: the flat law at its wall, the rounding lemma,
and the carrier floor.
\end{definition}

\begin{definition}[interlacement graph]\label{def:interlace}
The \emph{interlacement graph} $G_P$ has vertex set $[m]=\{1,\dots,m\}$, with
$c\sim c'$ if and only if exactly one of the two occurrences of $c'$ lies between
the two occurrences of $c$ in the Gauss word. (The relation is symmetric.) We
write $\Ind(G_P)$ for the set of independent sets of $G_P$, including
$\varnothing$, and $N_{G_P}(S)$ for the set of vertices adjacent to some element
of $S$.
\end{definition}

\begin{definition}[oriented smoothing]\label{def:smoothing}
For $S\in\Ind(G_P)$, the \emph{oriented smoothing of $P$ along $S$} replaces, at
each double point of $S$, the two transversally crossing arcs by the two arcs
that respect the orientation of $P$ and do not cross. The result is a disjoint
union of closed oriented curves, called the \emph{carriers} of $S$.
\end{definition}

\begin{lemma}[the carriers of a support]\label{lem:carriers}
Let $P$ be diagrammatic (Definition~\ref{def:diagrammatic}) and
$S\in\Ind(G_P)$ --- genericity is not needed here and an earlier revision
assumed it; what the argument reads is the Gauss word. Let $\Gamma$ be the
\emph{traversal circle}: an abstract oriented circle
together with the traversal map $\Gamma\to\mathbb R^2$ that runs once around
$P$, on which each double point of $P$ has exactly two preimages, so that
$\Gamma$ carries $2m$ marked points. (No coordinate is placed on $\Gamma$; in
particular the crossing-free case $m=0$ is allowed, with no marked points.) Then:
\begin{enumerate}
\item[(i)] the oriented smoothing of $P$ along $S$ has exactly $|S|+1$ carriers;
\item[(ii)] the assignment of traversal points to carriers is \emph{non-crossing}:
there are no four points $u_1,u_2,u_3,u_4$ of $\Gamma$ in this cyclic order,
none of them a preimage of an element of $S$, with $u_1,u_3$ on one carrier and
$u_2,u_4$ on a different carrier;
\item[(iii)] if $c\in[m]\setminus S$ is non-adjacent in $G_P$ to every element of
$S$, then both occurrences of $c$ lie on the same carrier;
\item[(iv)] if $H$ is a connected component of the induced graph
$G_P\bigl[[m]\setminus(S\cup N_{G_P}(S))\bigr]$, then there is exactly one
carrier on which both occurrences of every $c\in H$ lie.
\end{enumerate}
\end{lemma}
\begin{proof}
Induct on $|S|$, proving (i), (ii) and (iii) together; the induction is on the
conjunction, and (iii) is what makes the step go through.

\emph{Base $S=\varnothing$.} There is one carrier, namely $P$ itself: (i) and
(iii) are immediate and (ii) is vacuous, all four points lying on that one
carrier.

\emph{Step.} Let $|S|\geq1$, fix $c\in S$ and put $S'=S\setminus\{c\}$, so
$S'\in\Ind(G_P)$ and $|S'|=|S|-1$. Since $S$ is independent, $c$ is non-adjacent
to every element of $S'$, so the inductive hypothesis (iii) applies to $S'$ and
$c$: both occurrences of $c$ lie on one and the same carrier $L$ of $S'$. Write $x_1,x_2\in\Gamma$ for its two preimages; they cut $\Gamma$ into two
arcs $\alpha,\beta$.

(i) Oriented smoothing at $c$ replaces the two crossing arcs at that double point
by the two non-crossing arcs respecting the orientation, which severs $L$ at
$x_1$ and $x_2$ and reconnects it as the two closed curves carried by
$L\cap\alpha$ and $L\cap\beta$; every other carrier of $S'$ is untouched, since
$c$ meets none of them. So the number of carriers rises by exactly one, from
$|S'|+1$ to $|S|+1$.

(ii) The carriers of $S$ are those of $S'$ with the single block $L$ replaced by
the two blocks $L\cap\alpha$ and $L\cap\beta$. Suppose $u_1,u_2,u_3,u_4$ in cyclic
order violate (ii) for $S$, with $u_1,u_3$ on a carrier $A$ and $u_2,u_4$ on a
carrier $B\neq A$, none of the $u_k$ an occurrence of an element of $S$. If
neither $A$ nor $B$ is one of the two new blocks, the same four points violate
(ii) for $S'$, contradicting the inductive hypothesis. If exactly one of them, say
$A$, is a new block, then $A\subseteq L$ and $B$ is a carrier of $S'$ distinct
from $L$, so the four points again violate (ii) for $S'$ with $L$ in place of $A$.
If both are new blocks, then $\{A,B\}=\{L\cap\alpha,\,L\cap\beta\}$, so
$u_1,u_3\in\alpha$ and $u_2,u_4\in\beta$ (or the same with $\alpha,\beta$
exchanged). None of the $u_k$ equals $x_1$ or $x_2$, so each of the four cyclic
transitions $u_1\to u_2\to u_3\to u_4\to u_1$ passes from $\alpha$ to $\beta$ or
back, and every such passage contains one of the two points $x_1,x_2$ in its
interior. Four disjoint open arcs would then each contain a point of the
two-element set $\{x_1,x_2\}$, which is impossible

(iii) Let $c''\in[m]\setminus S$ be non-adjacent to every element of $S$, hence to
every element of $S'$. By the inductive hypothesis both occurrences of $c''$ lie
on one carrier $L''$ of $S'$. If $L''\neq L$ then $L''$ is a carrier of $S$ as
well and the claim holds. If $L''=L$, use that $c''$ is non-adjacent to $c$:
exactly one occurrence of $c$ lying between the two occurrences of $c''$ is what
adjacency means, so non-adjacency puts both occurrences of $c''$ strictly inside
the same arc, $\alpha$ or $\beta$. Hence both lie on the same one of
$L\cap\alpha$, $L\cap\beta$.

(iv) Put $U=[m]\setminus(S\cup N_{G_P}(S))$. For $c\in H$, clause~(iii) puts
both occurrences of $c$ on one carrier, which we call $L(c)$. If $c,c'\in H$
are adjacent in $G_P[U]$, exactly one occurrence of $c'$ lies between the two
occurrences of $c$. Thus their four occurrences alternate on $\Gamma$, and none
is an occurrence of an element of $S$. If $L(c)\neq L(c')$, those four points
violate clause~(ii). Hence $L(c)=L(c')$. The assignment $c\mapsto L(c)$ is
therefore constant on every edge of the connected graph $H$, hence constant on
$H$.

It is one carrier, not two, because it is a function. Uniqueness is immediate:
the carriers partition $\Gamma$, so distinct carriers are disjoint \emph{as
subsets of the traversal circle} --- not as subsets of the plane, where two
distinct carriers can meet only at \emph{dominated} crossings, those of
$N_{G_P}(S)$. At a crossing outside $S\cup N_{G_P}(S)$ both visits lie on one
carrier by clause~(iii), so distinct carriers cannot meet there, and the
crossings of $S$ are smoothed. An earlier revision said distinct carriers meet
at every shared double point \emph{outside} $S\cup N_{G_P}(S)$, which names
exactly the set where they cannot.
\end{proof}

\begin{lemma}[smoothing preserves the induced cyclic order]\label{lem:carrierword}
Let $C$ be a closed curve with a traversal circle $\Gamma_C$ and finitely many
transverse double points, and let $S$ be a set of them no two of which
interlace. Then each closed curve produced by smoothing $C$ along $S$ traverses
the marked points of $\Gamma_C$ lying on it in the cyclic order they have on
$\Gamma_C$.

In particular, taking $C=P$ generic and $S\in\Ind(G_P)$, each carrier of $S$
traverses the marked points lying on it in the order induced from $\Gamma$.
\end{lemma}
\begin{proof}
Induction on $|S|$. If $S=\varnothing$ the only curve is $C$ itself, traversing
$\Gamma_C$, and there is nothing to prove.

Let $|S|\geq1$ and pick $c\in S$, with traversal preimages $c_1,c_2$. Smoothing
$c$ alone joins the incoming branch at $c_1$ to the outgoing branch at $c_2$ and
the incoming at $c_2$ to the outgoing at $c_1$, so the two resulting closed
curves traverse exactly the two arcs $\Gamma_1=(c_1,c_2)$ and
$\Gamma_2=(c_2,c_1)$ of $\Gamma_C$, each in the order induced from $\Gamma_C$;
no marked point changes arcs and none is visited twice.

Every other $e\in S$ is nonadjacent to $c$, so its two preimages do not separate
$c_1$ from $c_2$: both lie in $\Gamma_1$ or both in $\Gamma_2$. Hence
$S\setminus\{c\}$ is partitioned into two sets, one carried by each of the two
closed curves, no two members of either interlacing; and the remaining
smoothings are performed inside those two curves separately. Each is a closed
curve with a traversal circle and finitely many transverse double points, which
is the hypothesis of this lemma, so the induction hypothesis applies to each and
gives the induced order there; composing with the previous paragraph gives it on
$\Gamma_C$.

The statement was restated at this generality because the induction applies it
to curves obtained by smoothing, which are not polygons; an earlier revision
quantified it over generic polygons only, and the peer was right that its own
proof then stepped outside its domain.
\end{proof}

\begin{definition}[corner, turn sign, and $\wind$]\label{def:wind}
Let $P$ be \emph{generic} and $S\in\Ind(G_P)$ --- the binders under which the
definition is read, and which an earlier revision left to the ambient
subsection, whose standing hypothesis is only that $P$ is diagrammatic. A
\emph{corner} of a carrier of $S$ is either a vertex $p_i$ of $P$ traversed by
it or a smoothing site of $S$ traversed by it. At each corner the carrier turns
\emph{left} or \emph{right} according to the sign of the determinant of the
incoming and outgoing directions, which under genericity is nonzero: at a
vertex by (G1), unconditional and hence relevant, at a smoothing site by (G5),
active at a crossing. On a merely diagrammatic polygon the vertex clause is
false --- $((0,0),(1,0),(2,0),(0,1))$ is diagrammatic and its turn at $p_2$ is
exactly zero --- which is why the generic binder is part of this definition and
not decoration. A carrier is \emph{uniform} if all its corners turn the
same way and \emph{mixed} otherwise. Its \emph{weight} is
\[
   \mathrm{wt}(L)=\begin{cases}
     +1, & L \text{ uniform, all turns right},\\
     (-1)^{c(L)}, & L \text{ uniform, all turns left, } c(L)=\#\text{corners},\\
     0, & L \text{ mixed},
   \end{cases}
\]
and $\wind(S)=\prod_{L}\mathrm{wt}(L)$, the product over the $|S|+1$ carriers of
$S$. In particular $\wind(S)\neq0$ forces every carrier of $S$ to be
uniform: $\wind(S)$ is a product of integers, a product of integers is
nonzero only if every factor is, and the weight of a mixed carrier is $0$.
\end{definition}

\begin{definition}[undominated set and residual pieces]\label{def:pieces}
For $S\in\Ind(G_P)$ put $U(S)=[m]\setminus\bigl(S\cup N_{G_P}(S)\bigr)$. The
\emph{residual pieces} of $S$ are the connected components
$H\in\pi_0\bigl(G_P[U(S)]\bigr)$ of the induced subgraph on $U(S)$.
\end{definition}

\subsection{The readout}

\begin{definition}[record, and record isomorphism]\label{def:record}
A \emph{record} consists of an oriented circle $\Gamma$; a finite subset
$V\subset\Gamma$ of \emph{marked points}; a partition of $V$ into two-element
sets, the \emph{double points}; for each double point a designation of one of
its two points as the \emph{over} position and the other as the \emph{under}
position; and for each double point a \emph{sign} in $\{\pm1\}$. The record of a
link diagram is the one obtained by taking $\Gamma$ to be its traversal circle,
$V$ the preimages of its crossings, the partition the crossing correspondence,
and the over/under and sign data those of the diagram.

A \emph{record isomorphism} from $(\Gamma,V,\dots)$ to $(\Gamma',V',\dots)$ is a
bijection $\varphi:V\to V'$ that
\begin{enumerate}
\item[(a)] preserves the cyclic order: for all $v_1,v_2,v_3\in V$ occurring in
that cyclic order on $\Gamma$, the points $\varphi v_1,\varphi v_2,\varphi v_3$
occur in that cyclic order on $\Gamma'$;
\item[(b)] carries double points to double points;
\item[(c)] carries over positions to over positions and under to under;
\item[(d)] preserves signs.
\end{enumerate}
Condition (a) says exactly that $\varphi$ extends to an orientation-preserving
homeomorphism $\Gamma\to\Gamma'$ carrying $V$ onto $V'$, and any two such
extensions are isotopic through such homeomorphisms: an orientation-preserving
circle homeomorphism is determined up to isotopy by the cyclic-order-preserving
bijection it induces on a finite subset, the complementary arcs being carried to
the complementary arcs in the induced order and any two orientation-preserving
homeomorphisms of a closed arc rel endpoints being isotopic. Records are
\emph{isomorphic} when such a $\varphi$ exists; the relation is an equivalence.
\end{definition}

\begin{definition}[HOMFLY--PT normalization]\label{def:homfly}
$P_H$ is normalized by the first two identities below; the third is their
consequence, obtained by applying the skein relation at a crossing between $L$
and a split unknot, and is displayed because it fixes the reader's convention
for a split component:
\begin{equation}
P(\bigcirc)=1,\qquad
aP(L_+)-a^{-1}P(L_-)=zP(L_0),\qquad
P(L\sqcup\bigcirc)=\frac{a-a^{-1}}{z}\,P(L).
\label{eq:homfly}
\end{equation}
\end{definition}

\begin{definition}[piece diagram, positivity, writhe]\label{def:piecediagram}
For a diagrammatic parent $P$, the diagram of a residual piece $H$ is the parent
curve $P$ with every double
point outside $H$ erased --- the two strands drawn as passing without
interaction --- and every double point of $H$ resolved by the divide convention:
the branch whose direction $u_{\mathrm{over}}$ satisfies
$\det(u_{\mathrm{over}},u_{\mathrm{under}})>0$ passes over. Under this convention
every crossing of the diagram is positive, so its writhe is
$w(H)=|H|$, the number of double points of $H$. We write $P_H(a,z)$ for the
HOMFLY--PT polynomial of the link so presented, normalized as in
Definition~\ref{def:homfly}.

\emph{What ``erased'' means, stated once, and where the statement comes from.}
The evaluator that defines $X_1$ builds a piece's datum as the parent Gauss word
restricted to the crossings of $H$, in the parent's cyclic order, together with
the parent's per-crossing counter-clockwise half-edge order restricted to the
surviving half-edges and re-expressed in the restricted word's arc labels
(\texttt{sub\_word\_and\_rot},
\texttt{engines/polygon\_motivic/ruling\_dp\_statesum\_fin.py:28} of the shipped
code). So the datum is a word \emph{together with a rotation system}, not a bare
word, and the definition above is to be read as naming that pair.

To erase a double point is to omit it from that datum. The combinatorial shadow
of the datum --- the word with its over/under and signs and nothing else --- is
Definition~\ref{def:record}.
\end{definition}

\begin{lemma}[a residual piece is carried by an actual closed curve]
\label{lem:piececurve}
Let $P$ be diagrammatic (Definition~\ref{def:diagrammatic}),
$S\in\Ind(G_P)$, and let $H$ be a residual piece of $S$. Let $C_H$ be the closed
plane curve that the proof's own route constructs, and nothing else: smooth
every crossing of $S$ --- of $S$ \emph{only} --- so that the curve falls into
the $|S|+1$ carriers of $S$; take the one carrier that carries $H$, which is
unique by Lemma~\ref{lem:carriers}(iv); and apply to it Step 5's iteration,
smoothing at each step one double point outside $H$ and keeping the daughter
curve that carries $H$, which Step 5 shows is well defined. Then $C_H$ is a
closed plane curve whose double points are exactly the crossings of $H$ and
whose Gauss word is the parent word restricted to $H$ in the parent's cyclic
order.

This binding is the construction and not a description of its result. The
route above smooths nothing but $S$, so $H$ lies on one carrier before the
iteration begins, and each subsequent step preserves that. In particular that restricted word is realizable, and the datum of
Definition~\ref{def:piecediagram} is the datum of $C_H$.
\end{lemma}
\begin{proof}
\emph{Step 1: a dominated crossing has its two visits on different carriers.}
Let $c\notin S$ interlace some $s\in S$. Smoothing $s$ alone splits the curve
into two closed curves traversing the two arcs $(s_1,s_2)$ and $(s_2,s_1)$ of
$\Gamma$, and $c$ interlacing $s$ means one visit of $c$ lies in each. Smoothing
the remaining elements of $S$ only ever splits: the carriers number $|S|+1$ by
Lemma~\ref{lem:carriers}(i), which is the count reached when every one of the
$|S|$ smoothings increases the component count, so none of them merges two
components. Hence the two visits of $c$ lie on different carriers.

\emph{Step 2: the double points of a carrier are exactly the undominated
crossings it carries.} A double point of a carrier $L$ is a crossing outside $S$
both of whose visits lie on $L$. Step 1 excludes the dominated ones, and
Lemma~\ref{lem:carriers}(iii) puts both visits of each undominated one on a
single carrier.

\emph{Step 3: a carrier's Gauss word is realizable and is a restriction.} A
carrier is a closed plane curve with transverse double points, so its Gauss word
is realizable, being read off an actual curve. By Step 2 its double points are
the undominated crossings it carries, and by Lemma~\ref{lem:carrierword} it
traverses them in the cyclic order they have on $\Gamma$. So the carrier's word
is the parent word restricted to those crossings.

\emph{Step 4: a residual piece interlaces nothing else in that word.} $H$ is a
connected component of $G_P[U(S)]$, so an undominated crossing interlacing a
member of $H$ lies in $H$; and $H$ lies on one carrier by
Lemma~\ref{lem:carriers}(iv).

\emph{Step 5: restricting to such a set is realizable, constructively.} Let $C$
be a closed plane curve and $K$ a nonempty set of its double points such that no
member of $K$ interlaces a double point outside $K$ \emph{and $K$ induces a
connected interlacement subgraph}. If $C$ has a double point $d$ outside $K$,
smooth it. The curve becomes two closed curves, traversing the two arcs of $d$.
Since $d$ interlaces no member of $K$, each member of $K$ has both visits in one
of those two arcs; and two members lying in different arcs have their visits in
disjoint arcs, so they do not interlace. Connectedness of $K$ therefore forces
all of $K$ into one arc, hence onto one of the two curves. Keep that one and
discard the other, which carries no member of $K$; the retained visits keep
their cyclic order by Lemma~\ref{lem:carrierword} applied to the single smoothed
chord. The hypotheses persist: $K$ is unchanged, so it is still connected, and
it still interlaces no surviving double point outside it. Each step removes at
least one double point outside $K$, so after finitely many steps the curve
$C_H$ has double points exactly $K$ and Gauss word the restriction.

The connectedness hypothesis is not decoration. Without it the step is false,
and the peer's exact witness is the word $d\,a\,a\,d\,b\,b$ with
$K=\{a,b\}$: neither $a$ nor $b$ interlaces $d$, yet $a$ lies inside the arc
$(d_1,d_2)$ and $b$ outside it, so $K$ is not on one side. A residual piece
supplies the hypothesis by Step 4, being a connected component.

Applying Step 5 to the carrier of $H$ with $K=H$, which Steps 3 and 4
licence --- Step 4 supplying both the non-interlacing and the connectedness
hypotheses --- performs exactly the statement's third operation and gives
$C_H$. Its rotation system at each surviving crossing is the carrier's,
which is the parent's, since smoothing at other points does not disturb the
half-edge order at a survivor; that is the second half of the datum.
\end{proof}

\begin{remark}[what the piece datum presents, and a reading that fails]
\label{rem:piecedatum}
An earlier revision of Definition~\ref{def:piecediagram} defined the diagram by
its record alone. That is not the evaluator's datum and it is not safe:
restricting a realizable Gauss word to an arbitrary label set can produce a word
no plane curve realizes --- the peer's witness restricts $123123$ to two of its
three crossings and gets $2323$, in which each chord interlaces exactly one
other, against the realizable-word parity used in the proof of
Corollary~\ref{cor:evendominated}. What rescues the
definition is not a reading but the lemma above: a residual piece is not an
arbitrary label set. The record shadow is what Lemma~\ref{lem:pieceintrinsic}
compares.

Definition~\ref{def:piecediagram} names a word together with a rotation system
and calls it the diagram of a piece. Lemma~\ref{lem:piececurve} is what makes
that a link diagram rather than a formal datum: the pair is the datum of an
actual closed plane curve, so $P_H$ is a classical invariant and no
realizability question is left open. The fact is recorded here, after the lemma,
rather than inside the definition, which cannot draw on a result printed after
it.
\end{remark}

\begin{corollary}[undominated crossings meet an even number of dominated ones]
\label{cor:evendominated}
Let $P$ be generic and $S\in\Ind(G_P)$. Every undominated crossing interlaces an
even number of dominated crossings.
\end{corollary}
\begin{proof}
Let $h$ be undominated, on the carrier $L$. The word of $L$ is realizable
(Step 3 above), so $h$ interlaces an even number of chords of that word, which
by Step 2 are the undominated crossings on $L$; and $h$ interlaces no
undominated crossing off $L$, by the non-crossing property
Lemma~\ref{lem:carriers}(ii). So $h$ interlaces an even number of undominated
crossings in $\Gamma$'s word. That word is realizable too, so $h$ interlaces an
even number of chords altogether, and $h$ interlaces no member of $S$, being
undominated. Subtracting leaves an even number of dominated ones.
\end{proof}
\begin{proof}[Second proof, the peer's]
This argument is not this author's and is printed with that attribution. For any
independent $T$, the set $N(T)$ is exactly the set of crossings whose two visits
lie on \emph{different} carriers of $T$, by Steps 1 and 2 above. Distinct closed
curves in the plane meet in an even number of transverse points, so summing over
the pairs of carriers, $|N(T)|$ is even. Now $h$ is undominated, so $S$,
$\{h\}$ and $S\cup\{h\}$ are all independent, and
\[
   N\bigl(S\cup\{h\}\bigr)=N(S)\cup N(h).
\]
Taking cardinalities modulo two and using that all three of
$|N(S\cup\{h\})|$, $|N(S)|$ and $|N(h)|$ are even gives $|N(S)\cap N(h)|$ even,
and that intersection is exactly the set of dominated crossings interlacing $h$.

The peer's caveat is worth keeping with the proof: this parity fact is necessary
for realizability but not equivalent to it, so it does not replace
Lemma~\ref{lem:piececurve}, whose carrier construction remains load-bearing.
\end{proof}


\begin{definition}[rotation number, computably]\label{def:rot}
Let $L$ be a closed polygon through corner points $q_0,\dots,q_{c-1}$, lying
in the regular locus $\mathcal R_c$ of Definition~\ref{def:regular}(B) ---
nonzero edges, no consecutive pair doubling back; equivalently, every
principal turn of Definition~\ref{def:regular}(A) exists --- and put
$\delta_i=q_{i+1}-q_i$. Choose $r\in\mathbb{R}^2$ with $\det(r,\delta_i)\neq0$
for all $i$; such $r$ exists because the $\delta_i$ are finitely many. Then
\[
   \rot(L)=\sum_{i}\epsilon_i,\qquad
   \epsilon_i=\begin{cases}
     +1,&\det(\delta_{i-1},\delta_i)>0,\ \det(\delta_{i-1},r)>0,\ \det(r,\delta_i)>0,\\
     -1,&\det(\delta_{i-1},\delta_i)<0,\ \det(\delta_{i-1},r)<0,\ \det(r,\delta_i)<0,\\
     0,&\text{otherwise.}
   \end{cases}
\]
All comparisons are exact sign tests of determinants: no angles and no
tolerances occur.

\emph{Why regularity is part of the definition.} The sum is independent of
the admissible $r$, and regularity is what makes it so; a vanishing turn is
harmless. First delete flat corners: if $\delta_i=\lambda\delta_{i-1}$ with
$\lambda>0$ then $\det(\delta_{i-1},\delta_i)=0$, so $\epsilon_i=0$ for every
$r$, and deleting $q_i$ merges the two edges into one of the same direction,
changing no other $\epsilon_j$, since every determinant a surviving term
reads is altered only by a positive scale. The deletion lands in
$\mathcal R_{c-1}$, and no polygon of any $\mathcal R_c$ is all-flat --- a
closed polygon with every turn flat would traverse one direction forever and
could not close --- so finitely many deletions reach a polygon with every
$\det(\delta_{i-1},\delta_i)\neq0$ and the same sum. On that polygon,
$\epsilon_i$ depends on $r$ only through $\det(\delta_{i-1},r)$ and
$\det(r,\delta_i)$, so it moves only as $r$ crosses the line of
$\delta_{i-1}$ or of $\delta_i$. Crossing the line of $\delta_j$ moves
exactly $\epsilon_j$ and $\epsilon_{j+1}$, and in each of the four sign cases
one gains what the other loses, so the sum is unchanged; if several edges
share that line they are pairwise nonconsecutive --- consecutive ones would
be flat or doubling back, and both are gone --- so their pairs are disjoint
and each is preserved separately. Regularity is not decoration: the polygon
$((-7,4),(-9,6),(9,-9),(3,3),(9,-9),(-7,-7),(-7,-6))$, whose consecutive
pair $(-6,12),(6,-12)$ doubles back, returns both $-1$ and $0$ at admissible
rays.

\emph{Direction loops and smooth curves.}
For a continuous loop $T\colon\mathbb R/\mathbb Z\to S^1$, a
\emph{tangent-angle lift} at a seam is a continuous
$\theta\colon[0,1]\to\mathbb R$ such that
$T(s)=(\cos\theta(s),\sin\theta(s))$. Put
\[
   \operatorname{tw}(T)=\frac{\theta(1)-\theta(0)}{2\pi}.
\]
For a closed $C^1$ regular oriented curve
$\gamma\colon\mathbb R/\mathbb Z\to\mathbb R^2$, put
\[
   T_\gamma=\frac{\gamma'}{|\gamma'|},
   \qquad
   \rot(\gamma)=\operatorname{tw}(T_\gamma).
\]
Lemma~\ref{lem:turnlift} constructs the lift and proves that these values do
not depend on its choices, on the seam, or on an orientation-preserving regular
reparametrisation. For either a polygon or a $C^1$ regular closed curve put
$R(L)=|\rot(L)|$.

\end{definition}

\begin{lemma}[tangent lifts and principal turns]\label{lem:turnlift}
Here $\rot$ is as in Definition~\ref{def:rot}.
\begin{enumerate}
\item[(i)] Every continuous direction loop has a tangent-angle lift.
The integer $\operatorname{tw}(T)$ is independent of the lift and seam,
is unchanged by an orientation-preserving reparametrisation, and is constant
under a continuous homotopy through direction loops. Consequently the rotation
of a closed $C^1$ regular curve is invariant under regular homotopy, and
reversing its orientation negates it.
\item[(ii)] If $L\in\mathcal R_c$ has principal turns $\tau_i$, then
\[
   2\pi\rot(L)=\sum_i\tau_i.
\]
Consequently rotation is constant along every path in $\mathcal R_c$, is
unchanged by a positive flat subdivision, and is negated by orientation
reversal. Every regular three-corner polygon has rotation $+1$ or $-1$
according to the common sign of its three turns.
\item[(iii)] Suppose a closed $C^1$ regular curve is obtained from $L$ by
replacing every corner by a regular arc whose compatible tangent-angle lift
has increment $\tau_i$. Its rotation equals $\rot(L)$.

More generally, let $b,c$ be regular oriented arcs having the same endpoints
and the same oriented tangent rays at both endpoints, and suppose replacing
$b$ by $c$ in a closed curve gives closed $C^1$ regular curves $F_b,F_c$.
For compatible tangent lifts with equal initial values, write their increments
as $\Delta_b,\Delta_c$. Then
\[
   \rot(F_c)-\rot(F_b)=\frac{\Delta_c-\Delta_b}{2\pi}.
\]
If $A_t$, $0\leq t\leq1$, is a supplied continuous path in
$\mathrm{GL}^+(2,\mathbb R)$ with $A_0=I$ and $A_1=A$, applying $A$ to both
arcs leaves $\Delta_c-\Delta_b$ unchanged.
\end{enumerate}
\end{lemma}
\begin{proof}
For~(i), uniform continuity gives
$0=s_0<\cdots<s_m=1$ such that
$T(s)\cdot T(s_{j-1})>0$ on every
$[s_{j-1},s_j]$. After choosing $\theta(0)$, continue it there by
\[
 \theta(s)=\theta(s_{j-1})+
 \operatorname{atan2}\!\left(
   \det(T(s_{j-1}),T(s)),\,
   T(s_{j-1})\cdot T(s)\right).
\]
The second argument is positive, so the added angle lies in
$(-\pi/2,\pi/2)$; the formulas match at the subdivision points and construct a
continuous lift without a covering-space theorem. Since $T(1)=T(0)$, its
endpoint increment belongs to $2\pi\mathbb Z$. Two lifts differ continuously
by a value in $2\pi\mathbb Z$, hence by one constant.

If the increment is $2\pi k$, extend the lift by
$\theta(s+n)=\theta(s)+2\pi nk$. Every interval of length one then has the same
increment, proving seam independence. An orientation-preserving
reparametrisation of the circle has an increasing representative
$\phi\colon\mathbb R\to\mathbb R$ with
$\phi(s+1)=\phi(s)+1$; the lift $\theta\circ\phi$ has the same increment.

For a homotopy $T(s,t)$ fix $t_0$. Uniform continuity gives a neighbourhood of
$t_0$ on which
$T(s,t)\cdot T(s,t_0)>0$ for every $s$. There is therefore
a unique continuous relative angle
\[
 \alpha(s,t)=\operatorname{atan2}\!\left(
 \det(T(s,t_0),T(s,t)),\,
 T(s,t_0)\cdot T(s,t)\right)
 \in(-\pi/2,\pi/2).
\]
If $\theta_0$ lifts $T(\,\cdot\,,t_0)$, then
$\theta_0+\alpha(\,\cdot\,,t)$ lifts $T(\,\cdot\,,t)$.
Periodicity gives $\alpha(1,t)=\alpha(0,t)$, so the endpoint increment is
locally, hence globally, constant in $t$. A regular homotopy supplies such a
homotopy of normalised tangents. For the oppositely oriented curve,
$-T_\gamma(1-s)$ has lift $\theta(1-s)+\pi$, whose increment is the negative
of the original one.

For~(ii), give the admissible reference vector $r$ of
Definition~\ref{def:rot} an angle $\rho$, and let
$\beta_i\in(\rho,\rho+2\pi)$ represent the direction of the $i$th edge.
The three cases in that definition give exactly
\[
   \tau_i=\beta_i-\beta_{i-1}+2\pi\epsilon_i:
\]
a positive crossing of the reference ray contributes $+2\pi$, a negative
crossing contributes $-2\pi$, and otherwise no correction occurs.
Summing cyclically cancels the $\beta$ terms and proves the displayed identity.

Along a path in $\mathcal R_c$, every $\tau_i$ varies continuously in
$(-\pi,\pi)$; hence the displayed integer is constant. A positive flat
subdivision inserts one zero turn and changes no other turn. Reversal negates
and reverses the turn list.

For the three-corner claim, if the vertices are $A,B,C$, all three cyclic turn
determinants equal $\det(B-A,C-A)$. This determinant cannot vanish: otherwise
the three nonzero collinear edge scalars sum to zero, and a cyclic sign change
is a consecutive negative-multiple pair, contrary to regularity. Thus all
three turns have one sign. Their sum is nonzero and has magnitude below
$3\pi$; being $2\pi$ times an integer, it has rotation $+1$ or $-1$ according
to that sign.

For~(iii), concatenate constant lifts on the straight pieces with the
prescribed corner lifts. Their total increment is $\sum_i\tau_i$, so~(ii)
proves the rounding assertion.

For the replacement assertion, use the same lift on the unchanged
complementary arc. Shifting that complementary lift at a join changes neither
its increment nor the computation; subtraction cancels it and leaves
$\Delta_c-\Delta_b$.

Finally follow the tangent path of $c$ and then the tangent path of $b$
backwards. This is a closed direction loop of increment
$\Delta_c-\Delta_b$. The maps
\[
   \nu_t(z)=\frac{A_tz}{|A_tz|}
\]
give a homotopy of that direction loop. Clause~(i) keeps its increment
constant, proving the last assertion without invoking degree.
\end{proof}

\begin{definition}[the slot, the factor, and $X_1$]\label{def:X1}
Let $P$ be generic and $S\in\Ind(G_P)$. For a carrier $L$ of $S$ set
\[
  P_{S,L}=\prod_{H\text{ carried by }L}P_H,\qquad
  w_{S,L}=\sum_{H\text{ carried by }L}w(H),
\]
the products and sums taken over the residual pieces assigned to $L$ by
Lemma~\ref{lem:carriers}(iv), with the empty product $P_{S,L}=1$ and the empty
sum $w_{S,L}=0$ when no piece is carried by $L$. Here $P_H$ is the HOMFLY--PT
polynomial of the link that the piece $H$ presents, in the normalization of
Definition~\ref{def:homfly}; that such a polynomial exists at all is
Axiom~\ref{ax:homfly}, and this definition is where it is first read. The \emph{slot} of $L$ is the
integer $1-w_{S,L}-R(L)$, the \emph{factor} is
\[
   \Omega_1(S,L)=\bigl[a^{\,1-w_{S,L}-R(L)}z^{0}\bigr]P_{S,L}(a,z)
\]
--- the coefficient itself, taken as a value, with no sign gate and no
zero-gate applied --- and
\[
   X_1(P)\;=\;\sum_{S\in\Ind(G_P)}\wind(S)\prod_{L}\Omega_1(S,L),
\]
the inner product running over the $|S|+1$ carriers of $S$.
\end{definition}

\begin{proposition}[chambers and chamber invariance]\label{prop:chamberinv}
Fix $n\geq3$.
\begin{enumerate}
\item[(i)] The generic locus $\mathcal U_n$ is open in $(\mathbb R^2)^n$, and
every chamber --- every connected component of it --- is open and path connected.
\item[(ii)] $X_1$ is constant on each chamber.
\end{enumerate}
\end{proposition}
\begin{proof}
For clause (i), fix $P\in\mathcal U_n$. The argument must control the members that are
\emph{not} relevant at $P$ as well as those that are: openness of an activation
region says that a member active at $P$ stays active nearby, and says nothing
about one that is inactive at $P$ becoming active. An earlier revision inferred
the second from the first.

What controls both is the unconditional layer. Every unconditional member ---
every $\mathrm{G1}_i$ and every $\mathrm{G2}_{e,i}$ --- is relevant at every
polygon, so all of them are nonzero at $P$, they are finitely many, and they are
polynomials; let $N$ be a neighbourhood of $P$ on which each of them keeps its
sign. On $N$ the crossing status of every pair of remote edges is constant,
being the four-sign condition of Definition~\ref{def:guarded} in exactly those
members; hence the activation of every conditional member is constant on $N$,
and so is the set of members relevant at a point. Shrinking $N$ so that the
finitely many members relevant at $P$ --- now the members relevant everywhere on
$N$ --- also keep their signs, every point of $N$ is generic. So
$\mathcal U_n$ is open.

A connected component of an open subset of $\mathbb R^N$ is open, the space
being locally connected, and an open connected subset of $\mathbb R^N$ is path
connected: the set of points reachable from a fixed one by a path in it is open
and its complement in the component is open, and the component is connected. An
earlier revision asserted both facts inside the proof that consumes them,
rather than proving them.

For clause (ii), let $P^{0},P^{1}$ lie in one chamber and choose a path
$t\mapsto P(t)$, $t\in[0,1]$, between them inside it, as clause (i) allows.
Every statement below is proved along that path.

\emph{The combinatorics.} Each unconditional member $g\in\mathcal G$ is a
polynomial, so the composite $t\mapsto g(P(t))$ is continuous and, by genericity,
nowhere zero along the path. Its sign is therefore constant on the connected
interval $[0,1]$. The four-sign crossing test of
Definition~\ref{def:guarded} therefore fixes the set of crossing pairs
along the path, and hence a conditional member is active at one parameter
exactly when it is active at every parameter. Each $g$ in this now-fixed active
set is a polynomial relevant at every $P(t)$, so $t\mapsto g(P(t))$ is likewise
continuous and nowhere zero by genericity. Its sign is therefore constant as
well. So: the set of double points, indexed by the pairs of edges carrying them, is
constant; and no two of them collide, for the following reason, which splits in
two cases and uses (G3) in one and (G4) in the other --- an earlier revision
announced it as "the (G4) one and not the (G3) one", which is the half that is
true when the two double points share an edge. Two double points can collide in two ways, and an earlier revision admitted
only one of them, saying that four distinct edges through a point "is not a
coincidence of any member". It is one. \emph{If the four edges are distinct}, the collision puts four edges through a
single point, and there are two sub-cases. If some three of the four are
pairwise nonparallel, they cross pairwise at that point, so their
$\mathrm{G3}$ member is active, and it vanishes, three concurrent lines sharing
a projective point; being active it is relevant, so genericity forbids it. If
no such triple exists, then two of the four are parallel, and being parallel
lines through a common point they are the same line: two distinct edges lying
on one line, each containing that point in its relative interior, hence
overlapping in a segment. Then an endpoint of one lies on the other, an
\emph{unconditional} $\mathrm{G2}$ vanishes, and genericity forbids that as
well. Either way no such collision occurs along a path of generic polygons.
\emph{Otherwise the two double points share an edge}, and there the concurrency
member says nothing --- the two carrying edges may be parallel --- so the
argument is the (G4) one: let them lie on a common edge
$e$, carried by $f$ and $g$, in the order the path presents them. The member
that decides their order along $e$ is
$\mathrm{G4}\langle e;f,g\rangle$, the member-valued accessor of
Definition~\ref{def:guarded}, and the order itself is the sign of the oriented
value $\widehat{\mathrm{G4}}_{e;f,g}=(t_f-t_g)\det(d_f,d_e)\det(d_g,d_e)$,
which is that member times the orientation sign $\epsilon(f,g)$. The member is
active exactly when both $f$ and $g$ cross $e$, which is the case, so it is
relevant and hence nonzero of constant sign along the path; and since
$\epsilon(f,g)$ is a constant of the index data, the oriented value is of
constant sign too. Only the nonvanishing is used below, so the conclusion is
the same under either reading. A collision would be $t_f=t_g$, so no collision occurs and the
order along each edge is constant. When $f$ and $g$ share a vertex $M$ the same
member does the work, its vanishing being that of
$\mathrm{G2}_{e,M}\mathrm{G1}_M$. Both factors are unconditional members, so on
this path --- along which every polygon is \emph{generic} --- both are nonzero;
in particular $\mathrm{G1}_M\neq0$, and that is genericity and not regularity.
Membership of $\mathcal R_n$ (Definition~\ref{def:regular}(B)) excludes only zero
edges and a consecutive pair that doubles back, and a flat vertex with
$\mathrm{G1}_M=0$ is regular. So the product vanishes only if
$\mathrm{G2}_{e,M}$ does, which puts $M$ on the edge $e$ and is forbidden here.
An earlier revision of this sentence said "the turn is nonzero on a regular
polygon", which is false of $\mathcal R_n$. An earlier revision credited
the constancy of the active (G3) signs for all of this; a concurrency member is
inactive when its three edges do not pairwise cross, and then it says nothing,
while two crossings on one edge can still swap. Hence the cyclic sequence of crossing visits along
the traversal --- the Gauss word --- is constant, and so are $G_P$, $\Ind(G_P)$,
each $U(S)$, each family of residual pieces, and, by
Lemma~\ref{lem:carriers}(iv), the assignment of pieces to carriers. Each
carrier is determined by the Gauss word and the smoothing sites, so the carriers
correspond canonically along the path, with corners varying continuously.

\emph{The weights.} At a vertex corner the turn sign is $\sgn\mathrm{G1}_i$,
constant; at a smoothing corner of the crossing of $e,f$ it is the sign of
$\mathrm{G5}_{e,f}$ up to the fixed sign of the smoothing convention, also
constant. So each carrier is uniform on the whole path or mixed on the whole
path, its corner count is constant, and $\mathrm{wt}(L)$ and $\wind(S)$ are
constant (Definition~\ref{def:wind}).

\emph{The rotations.} Turn signs alone do \emph{not} determine a rotation
number --- the convex regular pentagon and the regular pentagram have the same
five positive turn signs and rotations $1$ and $2$ --- so this step is a path
argument, not a sign count. Along the path each carrier $L(t)$ has corners
varying continuously and no vanishing turn determinant, so each principal turn
$\tau_k(t)$ is continuous with values in $(-\pi,\pi)$, and
$\rot(L(t))=\frac1{2\pi}\sum_k\tau_k(t)$ by Lemma~\ref{lem:turnlift}(ii) is a
continuous integer-valued function of $t$ on the connected interval, hence
constant. So each $R(L)$ is constant.

\emph{The piece polynomials.} Fix two parameters $t_0,t_1$ and a residual
piece $H$, identified at the two parameters by the common Gauss word. By
Lemma~\ref{lem:piececurve}, each piece datum is realized by a connected closed
plane curve whose double points are exactly the crossings of $H$ and whose
cyclic word is the parent word restricted to $H$. Match each marked crossing--
carrying-edge incidence $(c,e)$ at $t_0$ with $(c,e)$ at $t_1$. This bijection
preserves cyclic order, and the two incidences bearing $c$ remain its crossing
pair. It preserves signs because every piece crossing is positive, and it
preserves over/under
positions because the divide convention chooses them from the sign of the
active $\mathrm{G5}$ member of the two carrying edges, which is constant along
the chamber. Thus it is an orientation-preserving record isomorphism
(Definition~\ref{def:record}). Lemma~\ref{lem:piececurve} retains exactly
the finite, distinct transverse crossings of $H$
and creates no self-intersection when it smooths locally; hence each realizing
curve has one component and no triple points. Axiom~\ref{ax:gausscode}
therefore says that the piece diagrams present the same oriented link. Hence
$P_H$ is constant by Axiom~\ref{ax:homfly}, while $w(H)=|H|$ is constant because
$H$ is.

Hence every $\Omega_1(S,L)=[a^{1-w_{S,L}-R(L)}z^0]P_{S,L}$ is constant along the
path, and so is the finite sum defining $X_1$. In particular
$X_1(P^0)=X_1(P^1)$.
\end{proof}

\begin{definition}[event, and its zero set]\label{def:event}
An \emph{event} is a continuous path $t\mapsto P(t)$ of polygons on $n$
vertices, $t\in(-\varepsilon,\varepsilon)$, such that $P(t)$ is generic for
every $t\neq0$ and $P(0)$ is not. Its \emph{zero set} is
\[
   Z=\bigl\{g\in\mathcal G:\ g(P(0))=0\ \text{ and }\ g\ \text{is relevant at }
   P(t)\ \text{for some }t\neq0\bigr\}.
\]
The \emph{two chambers of the event} are the chambers containing
$P\bigl((-\varepsilon,0)\bigr)$ and $P\bigl((0,\varepsilon)\bigr)$; each of
those two sets is connected and consists of generic polygons, so each does lie
in one chamber. We say the event is \emph{guarded relative
to $Z$}, and every appeal to simplicity below names its $Z$.

The event is \emph{transversal} if every member of $Z$ changes sign at $t=0$.
Every named event of the convention that lists them is required to be
transversal. Without that requirement the definition admits a tangency: the
path
\[
   p_1=(0,0),\quad p_2(t)=(2,-t^2),\quad p_3=(1,0),\quad p_4=(3,2)
\]
is generic for $t\neq0$ and has exactly the cusp zero bundle at $t=0$, yet
$P(t)=P(-t)$, so its two punctured sides are the same chamber and no crossing
is created. A tangency is not a wall crossing and every wall law would be
vacuous or false there.
\end{definition}

\begin{remark}[why the two chambers give two values]\label{rem:eventvalues}
$X_1$ takes one value on each of the two chambers of an event, by
Proposition~\ref{prop:chamberinv}(ii), which is what makes the jump across the event
a well-defined number and is the reason the definition names the chambers at
all. The fact is recorded here rather than in the definition, so that the
definition depends on nothing proved after it.
\end{remark}

\begin{lemma}[the guards outside the zero set]\label{lem:guardconst}
Let $t\mapsto P(t)$, $t\in(-\varepsilon,\varepsilon)$, be an event with zero set
$Z$, and let $g\in\mathcal G$ be relevant at $P(t)$ for some $t\neq0$ and not a
member of $Z$. Then $g(P(0))\neq0$, and there is
$\varepsilon'\in(0,\varepsilon]$ such that $g$ is nonzero and of constant sign
on the whole of $(-\varepsilon',\varepsilon')$.
\end{lemma}
\begin{proof}
By Definition~\ref{def:event}, $Z$ consists of the members that vanish at
$P(0)$ \emph{and} are relevant at some $t\neq0$. The hypothesis supplies the
second clause and denies membership, so the first clause fails: $g(P(0))\neq0$.
Since $g$ is a polynomial in the coordinates and $t\mapsto P(t)$ is continuous,
$g\circ P$ is continuous and nonzero at $0$, hence nonzero on some interval
$(-\varepsilon',\varepsilon')$, on which, being continuous and nowhere zero, it
has constant sign.
\end{proof}

\begin{remark}[why the conclusion is local and not global]\label{rem:guardlocal}
An earlier revision concluded that $g$ is nonzero on the \emph{whole} interval.
That does not follow and its own proof conceded as much: $g$ may be relevant on
one side of the event and inactive on the other, and an inactive conditional
member is under no constraint --- it may vanish, or vanish identically. What is
true, and what every consumer uses, is the local statement above. Since an
event is used only as a germ at $t=0$, shrinking $\varepsilon$ costs nothing,
and every appeal below is to the shrunken interval.
\end{remark}

\begin{remark}[why the zero set is defined and not imposed]\label{rem:zeroset}
An earlier form of this definition required \emph{every} member of
$\mathcal G\setminus Z$ to be nonzero throughout. That is too strong and makes
the intended domains empty: a conditional member that is inactive on both sides
of the event may vanish identically --- for instance a crossing-order predicate
$\mathrm{G4}_{e;f,g}$ for a pair $f,g$ that never crosses $e$ --- so no event
would satisfy it. It was also too weak in the other direction, since it did not
require the two punctured sides to be generic, and without that the values
$X_1(P_\pm)$ that every wall law compares need not exist.
Definition~\ref{def:event} asks only for genericity off the wall and reads the
zero set off the path; Lemma~\ref{lem:guardconst} then \emph{proves} the
constancy that used to be assumed.

The specification's Definition~3 declares an event simple when \emph{the}
predicate defining the named event vanishes and every other guarded predicate
keeps its sign. Read with $|Z|=1$ that has empty domain at some events used
below: at a flat event three members vanish together --- the turn
$\mathrm{G1}_j$ and the two incidences it forces --- so no such event has a
singleton zero set. The set-valued form above is what every appeal here uses.
\end{remark}

\begin{convention}[the named events]\label{conv:events}
Throughout, all three named events are required to be transversal in the
sense of Definition~\ref{def:event}, and for a polygon on $n$ vertices and an
index $j$:
\begin{itemize}
\item write $\mathcal T_j=\{\mathrm{G1}_j,\ \mathrm{G2}_{(j-1,j),\,j+1},\
\mathrm{G2}_{(j,j+1),\,j-1}\}$ for the \emph{turn trio at $j$}, and
\[
   \mathcal F_j=\bigl\{\,\mathrm{G4}\langle h;\,e_{j-1},e_j\rangle\ :\ h\
   \text{an edge remote to the vertex } p_j\,\bigr\}
\]
for the order predicates that compare the two crossings a single edge makes with
the two edges at $j$. The entries of $\mathcal F_j$ are \emph{members} of
$\mathcal G$ --- the member-valued accessor of Definition~\ref{def:guarded}
returns the member indexed by the reduced representatives --- as
Definition~\ref{def:event} requires of anything that can lie in a zero set. The
factorization is a statement about signs and is therefore written on the
oriented value:
$\widehat{\mathrm{G4}}_{h;e_{j-1},e_j}=-\mathrm{G2}_{h,j}\,\mathrm{G1}_j$,
where the orientation sign $\epsilon(e_{j-1},e_j)$ is $-1$ exactly when the
pair wraps, that is exactly when $j\equiv1$. Since $\epsilon$ is a nonzero
constant, the member and the oriented value vanish together and change sign
together, so the conclusion below may be read off either. On a
neighbourhood of a configuration with $\mathrm{G2}_{h,j}\neq0$ it is a nonzero
multiple of $\mathrm{G1}_j$: it vanishes with the trio and changes sign with it.
\item a \emph{simple flat event at $j$} is an event whose zero set is exactly
$\mathcal T_j$, at which $p_j$ lies strictly between $p_{j-1}$ and $p_{j+1}$ at
$t=0$. No member of $\mathcal F_j$ can be relevant at such an event, so nothing
is lost by requiring the zero set to be the trio alone: at $t=0$ the two edges
at $j$ lie on one line $\ell$ and meet only at $p_j$, which is interior to their
union. A remote edge $h$ is then in one of two positions. If $h$ is not
contained in $\ell$ it meets $\ell$ in at most one point, so it cannot meet both
legs in their relative interiors and cannot cross both. If $h$ \emph{is}
contained in $\ell$ it is parallel to both legs, and a parallel pair does not
cross --- the four-sign condition of Definition~\ref{def:guarded} fails, its
four $\mathrm{G2}$ members vanishing --- so it crosses neither. Either way no
member of $\mathcal F_j$ is active at $t=0$.

Activation on a punctured side is a germ condition and is settled in the
limit, not by the value at $t=0$ alone. Suppose $h$ crosses both legs at
parameters $t_n\neq0$ with $t_n\to0$; let $x_n$ and $y_n$ be the crossing
points on the two legs. The vertices are continuous in $t$ and the segments
are compact, so after passing to a subsequence $x_n\to x$ and $y_n\to y$, with
$x$ and $y$ in the closed edge $h(0)$ and in the respective closed legs at
$t=0$, all three of which lie where their $t=0$ positions are. Two positions of
$h(0)$ against the flat line $\ell$ remain, and both force $p_j\in h(0)$. If
$h(0)$ is not contained in $\ell$, it meets $\ell$ in at most one point; both
$x$ and $y$ lie in $h(0)\cap\ell$, so $x=y$, and that one point lies in both
closed legs, whose intersection is $\{p_j\}$. If $h(0)$ \emph{is} contained in
$\ell$, then it is a segment of $\ell$ meeting both closed legs, which lie on
opposite sides of $p_j$ along $\ell$, $p_j$ being interior to their union; a
connected set meeting both sides contains the point between them. Either way
$p_j\in h(0)$, so the unconditional member $\mathrm{G2}_{h,j}$ vanishes at
$t=0$, joins the zero set, and the simplicity requirement excludes it. A
previous repair of this paragraph argued instead from "$h$ meets at most one
leg's relative interior at $t=0$", which the contained-line position falsifies
--- legs $(-2,0)$--$(0,0)$--$(2,0)$ with $h(0)=[(-1,0),(1,0)]$ meet in both
relative interiors --- while the conclusion survives, $h(0)$ containing $p_j$
there. So on a
punctured neighbourhood of $0$ no member of $\mathcal F_j$ is relevant, which is
the condition Definition~\ref{def:event} imposes. An earlier revision argued the
first position only, writing "a straight edge $h$ meeting that line once", which
says nothing about an $h$ lying along it;
\item a \emph{simple cusp event at $j$} is an event with a zero set of that same
form --- $\mathcal T_j$ together with the relevant members of $\mathcal F_j$ ---
at which $p_j$ lies outside the segment $[p_{j-1},p_{j+1}]$ at $t=0$. The event
is \emph{threaded} when some member of $\mathcal F_j$ is relevant, that is when
some edge crosses both collapsing legs, and \emph{unthreaded} otherwise; both
are simple, and the laws are stated for both.

This is where an earlier revision of this document narrowed the domain. It
required the zero set to be exactly $\mathcal T_j$, which excludes every
threaded cusp. The witness, printed in Remark~\ref{rem:witnesses} in its
own indexing, has the cusp at $j=1$ on seven vertices: the trio vanishes as
$\tfrac25\to0\to-\tfrac25$ across the wall and the thread's member
$\mathrm{G4}_{4;1,7}$ as $-\tfrac{16}5\to0\to\tfrac{16}5$ --- oriented value
$\tfrac{16}5\to0\to-\tfrac{16}5$, the pair at $j=1$ wrapping --- the
thread crossing both legs on both sides --- while the cusp law and the main
theorem are both stated "threaded or not".
Phase 1.5 admitted those events under \emph{relative} simplicity; the clause
above is that condition printed. The enlarged zero set is still \emph{one}
condition: every added member is a nonzero multiple of $\mathrm{G1}_j$ near the
wall, so it vanishes and changes sign exactly when the turn does, which is the
flat bundle's pattern and is what the transversality requirement asks;
\item a \emph{simple vertex-on-edge event}, with moving vertex $M$ and remote
edge $(a,b)$, is an event at which $p_M$ lies in the relative interior of
$[p_a,p_b]$ at $t=0$, and whose zero set is \emph{branch-specific}:
\[
   Z_{\mathrm{bigon}}
   =\bigl\{\mathrm{G2}_{(a,b),\,M},\
     \mathrm{G4}\langle (a,b);\,e_{M-1},e_{M}\rangle\bigr\},
   \qquad
   Z_{\mathrm{sliding}}=\bigl\{\mathrm{G2}_{(a,b),\,M}\bigr\}.
\]
In the \emph{bigon} branch the two neighbours of $M$ lie strictly on one side of
the remote line, so on the two-crossing chamber both edges incident to $M$ meet
the remote edge; the order predicate comparing their crossing parameters along
$(a,b)$ is active there, and at the event the two parameters coincide at $p_M$,
so it vanishes. It is therefore relevant at some $t\neq0$ and vanishes at
$t=0$: by Definition~\ref{def:event} it belongs to $Z$, and a singleton zero
set would have empty domain, exactly as at a flat event.

In the \emph{sliding} branch the two neighbours lie strictly on opposite sides,
so on \emph{each} punctured side exactly one incident edge meets the remote
edge. That order predicate is then active on neither side; it vanishes at the
wall but is relevant at no $t\neq0$, so Definition~\ref{def:event} does not put
it in $Z$, and imposing it would empty the branch
\end{itemize}
\end{convention}

\begin{definition}[silent event]\label{def:silent}
An event is \emph{silent} if every member of its zero set $Z$ is a predicate
$\mathrm{G2}_{e,i}$ whose vanishing at $t=0$ places $p_i$ on the \emph{line} of
$e$ but not on the segment $e$. In particular no member of $Z$ is a turn (G1),
a crossing-order predicate (G4), a concurrency (G3) or a direction determinant
(G5).
\end{definition}

\begin{remark}[two clauses that were removed because they are empty]
\label{rem:silentclauses}
An earlier revision of Definition~\ref{def:silent} admitted two further kinds of
member of $Z$: a conditional (G3) whose common point at $t=0$ misses one of the
three segments, and a (G5) whose pair of edges has relative interiors not
meeting at $t=0$. Neither can occur.

The first is Lemma~\ref{lem:dictionary}(A)(a). For the second, split on whether
the two closed segments meet at $t=0$. If they do not, it is
Lemma~\ref{lem:dictionary}(A)(b). If they do, then since they are parallel
with disjoint relative interiors they meet at a point that is an endpoint of at
least one of them, say the vertex $p_i$ of $e$ lying on $f$; then
$\mathrm{G2}_{f,i}$ vanishes at $t=0$ \emph{with the vertex on the segment}, and
$\mathrm{G2}$ is unconditional, hence relevant at every $P(t)$, so that
predicate lies in $Z$ and the event is not silent, since
Definition~\ref{def:silent} asks of every member of $Z$ that its vanishing put
the vertex on the \emph{line} and not on the segment; no clause about its (G5)
would have made it so.

Removing the clauses does not change the class of silent events. What it costs
is that the proof of Lemma~\ref{lem:silence} can no longer say ``clause~(ii)
excludes that''; it argues instead from the dichotomy the relevance filter
provides, which is the honest form of the same step. The peer raised this; the
entry records it.
\end{remark}


\begin{lemma}[$X_1$ is unchanged across a silent event]\label{lem:silence}
Let $t\mapsto P(t)$ be a silent event. Then $X_1(P_+)=X_1(P_-)$.
\end{lemma}
\begin{proof}
By the proof of Proposition~\ref{prop:chamberinv}(ii), $X_1$ is determined by the
set of double points indexed by the pairs of edges carrying them, their order
along each edge, the turn signs at vertices and at smoothing sites, the
\emph{absolute rotation} $R(L)$ of each carrier, and the polynomial $P_H$ of each
residual piece. An earlier revision omitted $R(L)$ from this list,
although the slot of Definition~\ref{def:X1} reads it. It suffices to show each is
constant on the whole interval, $t=0$ included.

A double point of the pair $(e,f)$ exists exactly when the relative interiors of
$e$ and $f$ meet, and by continuity it can appear or disappear only through a
configuration in which an endpoint of one segment lies on the other segment ---
a $\mathrm{G2}$ predicate vanishing \emph{with the vertex on the segment}.
Definition~\ref{def:silent} excludes that --- a member of $Z$ may put $p_i$ on
the line of $e$ but never on the segment --- and every other relevant
$\mathrm{G2}$ is nonzero throughout by Lemma~\ref{lem:guardconst}.
So the set of double points is constant, and with it the set of active
conditional members. Nor do two double points on \emph{disjoint} pairs of edges collide: that puts
four edges through one point, and the two sub-cases are those of
Proposition~\ref{prop:chamberinv}(ii). If some three of the four are pairwise
nonparallel they cross pairwise there, so their $\mathrm{G3}$ member is active
and vanishes, and an active member is relevant, hence lies in $Z$ if it
vanishes at $t=0$ --- which Definition~\ref{def:silent} forbids, its zero sets
containing no (G3). Otherwise two of the four are collinear and overlap at the
point, so an unconditional $\mathrm{G2}$ vanishes there with the vertex on the
\emph{segment}, which Definition~\ref{def:silent} also forbids.
The collision would leave the Gauss word alone but would destroy the
diagrammatic hypothesis the isotopy step below reads.

Two double points on a common edge exchange order only by colliding, which puts
three segments through one common point. Fix the common edge $e$ and the two edges $f,g$ carrying the colliding double
points; both are remote to $e$, a crossing requiring remoteness. The predicate
that decides the order of those two crossings along $e$ is the member
$\mathrm{G4}\langle e;f,g\rangle$ --- the accessor, since $f$ and $g$ are
presented in the order the geometry gives and their representatives need not be
increasing --- and it is the right tool: it is \emph{active} exactly when $f$
and $g$ both cross $e$, which is the situation in hand, and it vanishes exactly
when the two crossing parameters coincide, which is the collision. Being active
it is relevant, so if it vanished at $t=0$ it would lie in the zero set $Z$;
and $Z$ contains no (G4) member, by Definition~\ref{def:silent} --- the zero
set holds members, so it is membership of the accessor's value that is
excluded, and the orientation sign is irrelevant to it. So it does not
vanish, and by Lemma~\ref{lem:guardconst} it keeps its sign: no collision.

An earlier revision reached instead for the concurrency member
$\mathrm{G3}_{e,f,g}$ and argued by cases on whether $f$ and $g$ are remote or
consecutive. That misses a configuration, and the peer's chart is exact: with
$e$ from $(-2,0)$ of direction $(4,0)$, $f_t$ from $(t,-1)$ of direction
$(0,2)$ and $g_t$ from $(-t,-2)$ of direction $(0,4)$, the triple is pairwise
remote, $f$ and $g$ are \emph{parallel} so $\mathrm{G3}$ is inactive and says
nothing, yet both cross $e$ and their crossings swap at $t=0$;
$\mathrm{G3}=-64t$ and $\mathrm{G4}_{e;f,g}=64t$, and it is the (G4) that is
active and in the zero set's way.

The two edges $f,g$ may also share a vertex $M$, and then the same predicate
does the work, being indexed for any pair remote to $e$. Its vanishing is then
the vanishing of $\mathrm{G2}_{e,M}\,\mathrm{G1}_M$. The second factor is a
turn, and it is nonzero here because $\mathrm{G1}_M$ is unconditional and every
polygon of a silent event's punctured sides is generic, with $\mathrm{G1}_M$
outside the zero set by Definition~\ref{def:silent} and so nonzero at $t=0$ as
well by Lemma~\ref{lem:guardconst} --- regularity alone would not give it, a
flat vertex being regular. So a collision would need $M$ on the line of $e$, and
on the \emph{segment}, since the collision point is the common crossing
--- a $\mathrm{G2}$ vanishing with the vertex on the segment, which
Definition~\ref{def:silent} excludes from $Z$ while $\mathrm{G2}$ is
unconditional and hence relevant.

So no collision occurs, the order along each edge is constant, and the Gauss
word with it.

The turn sign at a vertex is $\sgn\mathrm{G1}_i$, which is unconditional,
hence relevant, and lies in no $Z$ of a silent event, so it is nonzero of
constant sign by Lemma~\ref{lem:guardconst}. The
turn sign at a smoothing site of the crossing of $e$ and $f$ is the sign of
$\mathrm{G5}_{e,f}$ up to the fixed sign of the smoothing convention. That
predicate is not in $Z$, so by the same dichotomy either it is nonzero on a
neighbourhood, or $e$ and $f$ carry no double point for $t\neq0$ and hence no
smoothing site whose turn could change.

The absolute rotation $R(L)$ of a carrier is next, and it needs its own
argument: turn signs alone do not determine a rotation number, and link
isotopy cannot substitute for it, plane rotation not being a knot-type
invariant. Each carrier $L(t)$ has, by the no-birth and no-collision facts
above, a fixed finite cyclic corner set, each corner varying continuously in
$t$. Each corner's turn determinant --- the $\mathrm{G1}$ at a vertex, the
$\mathrm{G5}$ at a smoothing site --- is nonzero through the wall by the two
preceding paragraphs, $t=0$ included, so each principal turn $\tau_k(t)$ is
continuous with values in $(-\pi,\pi)$, and
$\rot(L(t))=\frac1{2\pi}\sum_k\tau_k(t)$ by Lemma~\ref{lem:turnlift}(ii) is a
continuous integer-valued function on the connected event interval, hence
constant. This is the path argument of Proposition~\ref{prop:chamberinv}(ii) run
through the wall, and it is printed rather than cited because a silent wall
lies inside no chamber. So each $R(L)$ is constant.

Finally, the common Gauss word identifies every residual piece $H$ on the two
sides. The edge-pair indexing is constant as proved above: a silent zero puts its vertex
off the remote segment, so no crossing changes carrying edge at the wall.
Lemma~\ref{lem:piececurve} realizes its datum on each side by a connected
closed plane curve whose double points are exactly the crossings of $H$. Match each marked crossing--carrying-edge
incidence $(c,e)$ on the first side with the same incidence $(c,e)$
on the second. The common restricted word makes this bijection cyclic-order
preserving; the two incidences bearing $c$ remain its crossing pair. Every sign
is positive, and the $\mathrm{G5}$ argument above shows that the divide
convention designates the same visit as over on the two sides. Thus the map is
an orientation-preserving record isomorphism
(Definition~\ref{def:record}). Lemma~\ref{lem:piececurve} retains exactly
the finite, distinct transverse crossings of $H$ from the two generic parents
and creates no self-intersection when it smooths locally; hence each realizing
curve has one component and no triple points. Axiom~\ref{ax:gausscode} says that the piece diagrams present the same
oriented link, so Axiom~\ref{ax:homfly} gives the same $P_H$; and
$w(H)=|H|$ is unchanged because $H$ is.
Every input to $X_1$ is therefore constant on $(-\varepsilon,\varepsilon)$, and
$X_1(P_+)=X_1(P_-)$.
\end{proof}

\begin{remark}[why this lemma is here]\label{rem:silencenew}
The induction that proves the main theorem transports a polygon along a path
and needs $\Delta=X_1-A$ to be unchanged wherever the path is not at one of the
four named walls. For $A$ that is a declared input (Axiom~\ref{ax:lawful}). For
$X_1$ it is two facts, not one: constancy inside a chamber
(Proposition~\ref{prop:chamberinv}(ii)) and invariance across the silent walls
between chambers (Lemma~\ref{lem:silence}). A silent wall \emph{is} a wall ---
a guarded predicate vanishes on it --- so chamber constancy does not cross it.
\end{remark}

\begin{remark}[the guarded list cannot be trimmed]\label{rem:chamberguards}
The proof of Proposition~\ref{prop:chamberinv}(ii) uses all five families, and an
implementation enforcing only some of them defines a coarser partition than
Definition~\ref{def:generic}(B), on which the proposition is false. (G1) fixes the
vertex turn signs, (G2) keeps vertices off edges, (G3) keeps three edges from
concurring, (G4) fixes the crossing order along each edge, (G5) fixes the
smoothing-site turn signs. Which pairs of edges cross, and hence which
conditional members are active, is decided by the (G2) signs alone. Dropping (G3) or (G4) permits two crossings to
exchange order along an edge with no guarded predicate changing sign; that
changes the Gauss word and can change $X_1$.
\end{remark}
